% Copyright (c) 2026 Geovana Neves. Licensed under CC BY 4.0 (https://creativecommons.org/licenses/by/4.0/). See LICENSE-AND-AUTHORSHIP.md.
%
% fsi/text/02-the-loop.tex

\section{The coupling loop}

\begin{frame}{The loop, one coupling call}
\begin{center}
\begin{adjustbox}{max width=\linewidth, max totalheight=0.54\textheight}
\begin{tikzpicture}[node distance=5mm]
  \node[fsibox, minimum width=3.3cm, minimum height=1.35cm] (fs)
    {\textbf{FlightStream}\\a steady iteration or\\a time step, on the\\current mesh};
  \node[fsibox, right=9mm of fs, minimum width=4.4cm, minimum height=1.35cm] (post)
    {\textbf{post-processing script} \texttt{fsi\_post.txt}\\[0.5mm]
     \texttt{UPDATE\_ALL\_SURFACE\_SECTIONS}\\
     \texttt{COMPUTE\_SURFACE\_SECTIONAL\_LOADS NEWTONS}\\
     \texttt{EXPORT\_SURFACE\_SECTIONAL\_LOADS}};
  \node[fsifile, right=9mm of post, minimum height=1.35cm, text width=3.7cm] (loads)
    {FS\_SurfaceSection\_Loads.txt\\[0.5mm]
     {\rmfamily\itshape per section: Fx, Fz, moment about the quarter chord}};
  \node[fsibox, below=15mm of loads, minimum width=3.9cm, minimum height=1.7cm,
        text width=3.8cm] (step)
    {\textbf{\texttt{pyfs-fsi} step}\\[0.5mm]
     loads moved to the elastic axis\\
     one beam solve per blade\\
     relaxed against the last call\\
     twist written as node translations};
  \node[fsifile, left=12mm of step, minimum height=1.35cm, text width=2.7cm] (disp)
    {FSIDisp.txt\\[0.5mm]
     {\rmfamily\itshape one translation per node, in import order}};
  \node[fsibox, below=15mm of fs, minimum width=3.3cm, minimum height=1.35cm] (morph)
    {the solver moves\\the blade nodes\\and deforms the mesh};
  \node[fsifile, right=6mm of step, yshift=4mm] (cfg) {config.json};
  \node[fsifile, below=2mm of cfg] (state) {state.json};

  \draw[fsiarrow] (fs) -- (post);
  \draw[fsiarrow] (post) -- (loads);
  \draw[fsiarrow] (loads) -- (step);
  \draw[fsiarrow] (step) -- (disp);
  \draw[fsiarrow] (disp) -- (morph);
  \draw[fsiarrow] (morph) -- node[fsinote, left] {next\\call} (fs);
  \draw[fsiarrow, CourseGold] (cfg.west) -- (cfg.west -| step.east);
  \draw[fsiarrow, CourseGold, {Stealth[length=2mm]}-{Stealth[length=2mm]}]
    (state.west) -- (state.west -| step.east);
\end{tikzpicture}
\end{adjustbox}
\end{center}
\vspace{-1mm}
{\small FlightStream's Aeroelastic Toolbox calls the executable once per
aeroelastic iteration of a steady row, or once per time step of an unsteady
one (\code{SET\_AEROELASTIC\_ITERATIONS 1}). Each call is stateless: what
persists is in \code{state.json}. Loads and displacements are in each blade's
own frame, or the wing's; the package never sees azimuth.\par}
\fsisource{FSI tutorial, ``The loop in one paragraph'', ``cli'';
  \code{cases/fsi\_workspace.py}, \code{wire\_workspace\_fsi};
  \code{fsi/\_\_init\_\_.py}}
\end{frame}

\begin{frame}{The structure: one Euler beam per blade, or along the wing}
\begin{columns}[T]
\begin{column}{0.52\textwidth}
\begin{itemize}\small
  \item One node per radial station, on the elastic axis; one member per
        bay, with the bay-averaged EI and GJ; the root station clamped.
  \item Unit elastic moduli, \(E = G = 1\): the section constants handed to
        PyNite are the configuration's EI and GJ, so no modulus is applied
        twice.
  \item Flap deflection \(w\) and elastic twist \(\theta\) carry mass
        (\(\mu\,l\) and \((I_1 + I_2)\,l\)); every other degree of freedom is
        condensed out of the eigenproblem exactly.
  \item Static solves are linear, or P-Delta whenever axial tension is
        present [6, 7].
\end{itemize}
\end{column}
\begin{column}{0.44\textwidth}
\begin{center}
\begin{adjustbox}{max width=\linewidth}
\begin{tikzpicture}[x=1cm,y=1cm]
  % the clamp
  \fill[pattern=north east lines, pattern color=SoftGray] (-0.35,-0.8) rectangle (0,0.8);
  \draw[thick, SoftGray] (0,-0.8) -- (0,0.8);
  % the beam and its stations
  \draw[very thick, CourseBlue] (0,0) -- (5.6,0);
  \foreach \x/\n in {0/0,1.4/1,2.8/2,4.2/3,5.6/4} {
    \fill[CourseBlue] (\x,0) circle (1.6pt);
  }
  \node[font=\tiny, below, text=SoftGray] at (0.7,0) {bay};
  \draw[SoftGray, decorate, decoration={brace, amplitude=3pt, mirror}]
    (1.4,-0.28) -- (2.8,-0.28) node[midway, below=2pt, font=\tiny] {EI, GJ averaged};
  % flap and twist at the tip
  \draw[-{Stealth[length=1.6mm]}, thick, CourseRed] (5.6,0) -- (5.6,0.9)
    node[above, font=\scriptsize] {\(w\), flap};
  \draw[-{Stealth[length=1.6mm]}, thick, CourseRed]
    (6.05,-0.35) arc[start angle=-60, end angle=200, x radius=0.25, y radius=0.45];
  \node[font=\scriptsize, CourseRed, right] at (6.2,0.1) {\(\theta\), twist};
  % the tension
  \draw[-{Stealth[length=1.6mm]}, thick, CourseGold] (3.0,0.45) -- (4.2,0.45)
    node[midway, above, font=\scriptsize, text=black] {\(\mu\,\Omega^2 r\)};
  \node[font=\tiny, text=SoftGray, align=center] at (-0.2,-1.15) {root\\clamped};
  \node[font=\tiny, text=SoftGray] at (5.6,-0.95) {tip};
\end{tikzpicture}
\end{adjustbox}
\end{center}
\begin{takeaway}
The stiffnesses in \code{config.json} are the beam's stiffnesses, typed or
generated, exactly as written.
\end{takeaway}
\end{column}
\end{columns}
\fsisource{FSI tutorial, ``beam''; \code{fsi/beam.py} module docstring}
\end{frame}

\begin{frame}{What rotation adds}
\begin{itemize}\small
  \item \textbf{Tension.} The distributed axial force \(\mu\,\Omega^2 r\)
        makes the tension \(N(r)\) emerge in the solver, and P-Delta turns it
        into bending stiffness, with no manual correction terms: the
        centrifugal stiffening of a rotating blade [8, 9, 10].
  \item \textbf{Propeller moment.} \(-\Omega^2 (I_1 - I_2)\sin\theta\cos\theta\)
        drives each section toward flat pitch. It depends on the total pitch,
        geometric plus elastic, so each call re-solves the beam a few times
        until the twist stops moving [9, 12].
  \item \textbf{In-plane softening.} The part of the flap deflection that
        lies in the plane of rotation feels the centrifugal field pulling it
        outward, a term \(\mu\,\Omega^2 \sin^2\beta\,w\), \(\beta\) the
        section's pitch, that lowers the stiffness; it is applied with the
        tension [8, 9].
  \item \textbf{Not modelled: the trapeze effect.} A declared limitation;
        omitting it overpredicts the elastic twist, which is conservative.
\end{itemize}
\fsisource{FSI tutorial, ``The loop in one paragraph'' and ``centrifugal'';
  \code{fsi/centrifugal.py} module docstring}
\end{frame}

\begin{frame}{A fixed wing: its own weight, and where the model stops}
\begin{itemize}\small
  \item \textbf{A wing that does not turn} has no centrifugal load. The
        structural solver applies the structure's \textbf{own weight}, gravity
        on its mass distribution, beside the aerodynamic loads, on a
        \code{steady} or an \code{unsteady} row; the loop is the same.
  \item \textbf{A blade held still} in a turning free stream
        (\code{qsteady\_rotor}, sector) is the opposite case: the aerodynamics
        see no motion, and the structural solver still applies the
        centrifugal loads at the row's rotation speed.
  \item \textbf{Quasi-static.} The mean deformation is exact within the beam
        idealisation; the once-per-revolution content is trustworthy only well
        below resonance, roughly \(n\,\Omega / \omega_n \le 0.3\) [12, 13].
        The convergence log prints this boundary in its header.
\end{itemize}
\begin{alertblock}{The formulas are transcribed, not re-verified}
\small The primary sources of the model have not been independently checked
against the implemented formulas; each formula lives in one small function
that cites its source, so a correction stays local.
\end{alertblock}
\fsidecided{\code{fsi/centrifugal.py}; the page ``FSI in a workspace''}
\end{frame}

\begin{frame}{Twist travels as three nodes per station}
\begin{center}
\begin{adjustbox}{max width=0.62\linewidth, max totalheight=0.36\textheight}
\begin{tikzpicture}[x=1cm,y=1cm]
  % undeformed section
  \draw[SoftGray, dashed] (0,0) ellipse (2.4 and 0.32);
  \draw[SoftGray] (-2.9,0) -- (2.9,0);
  \foreach \x in {-1.2,0,1.2} \fill[SoftGray] (\x,0) circle (2pt);
  % deformed section: translated by w and turned by theta
  \begin{scope}[shift={(0,1.15)}, rotate=9]
    \draw[CourseBlue, thick] (0,0) ellipse (2.4 and 0.32);
    \fill[CourseBlue] (-1.2,0) circle (2.2pt);
    \fill[CourseRed] (0,0) circle (2.2pt);
    \fill[CourseBlue] (1.2,0) circle (2.2pt);
  \end{scope}
  \draw[-{Stealth[length=1.6mm]}, CourseRed, thick] (0,0) -- (0,1.12)
    node[midway, left, font=\scriptsize] {\(w\)};
  \draw[-{Stealth[length=1.4mm]}, CourseBlue] (1.2,0) -- (1.2,1.33)
    node[midway, right, font=\scriptsize] {\(w + \theta d\)};
  \draw[-{Stealth[length=1.4mm]}, CourseBlue] (-1.2,0) -- (-1.2,0.97)
    node[midway, left, font=\scriptsize] {\(w - \theta d\)};
  \node[font=\tiny, text=SoftGray] at (1.2,-0.55) {toward LE};
  \node[font=\tiny, text=SoftGray] at (0,-0.55) {elastic axis};
  \node[font=\tiny, text=SoftGray] at (-1.2,-0.55) {toward TE};
  \draw[SoftGray, |-|] (0,-0.9) -- (1.2,-0.9) node[midway, below, font=\tiny] {\(d = 0.25\,c\)};
\end{tikzpicture}
\end{adjustbox}
\end{center}
\vspace{-1mm}
\begin{columns}[T]
\begin{column}{0.48\textwidth}
\begin{itemize}\footnotesize
  \item \code{FSIDisp.txt} carries translations only, so twist is encoded by
        three nodes per station, \(d = 0.25\,c\) apart
        (\code{node\_offset\_chord\_fraction}).
  \item Rigid section kinematics, \(\delta y = w + \theta\,d\): linear, so the
        inverse is exact to machine precision.
\end{itemize}
\end{column}
\begin{column}{0.48\textwidth}
\begin{itemize}\footnotesize
  \item One generator makes the node file, the ordering map and the
        \code{FSIDisp.txt} rows: the solver applies them strictly in import
        order, and a drift would corrupt the geometry in silence.
  \item One node file serves every blade, imported once per blade frame, in
        blade order.
\end{itemize}
\end{column}
\end{columns}
\fsisource{FSI tutorial, ``kinematics'' and ``nodes'';
  \code{fsi/nodes.py} module docstring; \code{fsi/config.py},
  \code{FsiConfig.node\_offset\_chord\_fraction}}
\end{frame}

\begin{frame}{The driver: four phases, keyed on the step counter}
\relax % a body opening with a brace group would be read as the frame subtitle
{\ftstablesize\renewcommand{\arraystretch}{1.2}
\begin{tabularx}{\linewidth}{@{}P{0.23\linewidth}R@{}}
\toprule
\textbf{Phase} & \textbf{Behaviour, and its key in \code{config.json} \code{phases}} \\
\midrule
1, wake development & zero displacements on the rigid blade;
  \code{wake\_development\_revolutions} (1.0) \\
2, averaged coupling & loads averaged over the window, relaxed updates
  \(d_{new} = d_{old} + \lambda (d_{calc} - d_{old})\);
  \code{averaging\_window\_revolutions} (1.0), \code{coupling\_relaxation} (0.4) \\
3, convergence watch & as 2, and converged when BOTH hold across consecutive
  revolutions: tip twist change below \code{tip\_twist\_tolerance\_deg} (0.05)
  and relative change of the integrated normal force below
  \code{thrust\_tolerance\_fraction} (0.02) \\
4, recording & instantaneous loads, \(\lambda = 1\) by design, twist recorded
  per step; \code{recording\_revolutions} (1.0) \\
\bottomrule
\end{tabularx}}
\par\vspace{1.5mm}
\begin{itemize}\scriptsize
  \item \textbf{Freshness is asserted}: each call must see an advancing solver
        iteration in the loads header; a repeat is refused, since it means more
        than one FSI iteration per time step.
  \item \textbf{Frozen mode}: a \code{fsi\_frozen\_displacements.txt} in the
        run folder is replayed verbatim at every call, with no loads and no solve.
\end{itemize}
\fsisource{FSI tutorial, ``driver and state''; \code{fsi/config.py},
  \code{PhaseSchedule} (defaults in brackets)}
\end{frame}

\begin{frame}{The structural nodes, and how the surface follows them}
\begin{columns}[T]
\begin{column}{0.52\textwidth}
\begin{itemize}\small
  \item \textbf{Inside the component.} The nodes lie on the camber line of
        each station, turned by its local twist, inside the blade or the
        wing; a node outside is refused at plan.
  \item \textbf{In the frame that turns with the blade}, so a displacement
        means the same thing at every position.
  \item \textbf{The surface by its boundary}: the aeroelastic surface is
        assigned by the boundary's own identifier in the mesh.
  \item \textbf{After the structural call}: a coupled point exports once the
        last call has moved the surface; a steady run waits for it.
\end{itemize}
\end{column}
\begin{column}{0.44\textwidth}
\begin{block}{A radial basis function carries the nodes to the panels}
\small The displacement known at the nodes is interpolated onto every
aerodynamic panel [14, 15, 16]. A beam is a \textbf{line} of nodes: a
compactly supported kernel reaches only the panels within its radius of that
line, and the rest of the surface barely moves. The coupled row uses the
\textbf{multiquadric} [14], whose support is global.
\end{block}
\end{column}
\end{columns}
\fsidecided{\code{fsi/nodes.py}; \code{cases/fsi\_workspace.py}; Guide 4,
  \code{aeroelastic\_rbf\_type}}
\end{frame}

\begin{frame}{Route B: the direct transform, a sanity check}
\begin{columns}[T]
\begin{column}{0.52\textwidth}
\begin{itemize}\small
  \item The same beam, solved on the solver's sectional loads, gives flap
        \(w(r)\) and twist \(\theta(r)\); \textbf{every mesh vertex is moved}
        about the elastic axis by that \(w\) and \(\theta\), and the deformed
        geometry is solved again, rigid. No interpolation, no in-run
        coupling.
  \item Repeated until the tip flap and twist stop moving, it is a
        partitioned static iteration: quasi-steady by construction.
  \item A zero deformation through the same chain reproduces the rigid
        answer, which is its own control.
\end{itemize}
\end{column}
\begin{column}{0.44\textwidth}
\begin{ftslimit}[What route B is for]
\small A simplified run, a check of the coupled route against an
independent one, and an FSI option where the coupled route is not
available. It is \textbf{not} the coupled route: it carries no dynamics and
no in-run deformation. The worked example of Part 7 was solved this way.
\end{ftslimit}
\end{column}
\end{columns}
\fsidecided{the worked example's route diagnostic [3]}
\end{frame}
